\documentclass[12pt]{article}
\usepackage[paperwidth=6in,paperheight=9in]{geometry}
\geometry{
  inner=15mm,
  outer=10mm,
  top=13mm,
  bottom=13mm,
  headheight=12pt,
  headsep=7mm,
  footskip=10mm,
  includehead=true,
  includefoot=true
}
\usepackage{fontspec}
\setmainfont{Libertinus Serif}
\usepackage{microtype}
\usepackage{setspace}
\usepackage{ragged2e}

\title{Bibliography \& Research Methodology}
\author{From \textit{Stolen Words}}
\date{}

\raggedbottom
\clubpenalty=150
\widowpenalty=150
\setlength{\parskip}{4pt plus 2pt minus 1pt}
\setlength{\parindent}{1em}

\begin{document}

\maketitle

\section*{Bibliography}

This section provides detailed source documentation for readers who wish to independently verify the claims made throughout this book. All verse comparisons, audio recordings, and textual changes can be verified through publicly available online sources.

\vspace{0.5cm}

\subsection*{Primary Textual Sources}

\textit{Bhagavad-gītā As It Is}, Original 1972 Macmillan Edition. A. C. Bhaktivedanta Swami Prabhupāda. New York: Macmillan, 1972.
\begin{itemize}
\item Online: asitis.com | prabhupadagita.com
\end{itemize}

\textit{Bhagavad-gītā As It Is}, Revised and Enlarged 1983 Edition. Bhaktivedanta Book Trust. Los Angeles, 1983.
\begin{itemize}
\item Online: vedabase.io
\end{itemize}

Vedabase Archives: Complete collection of Srila Prabhupāda's recorded lectures, class transcripts, letters, and conversations, 1966-1977.
\begin{itemize}
\item Online: vedabase.io | prabhupadavani.org | vanisource.org
\end{itemize}

"Bhagavad Gita Changes — Complete List." \textit{ISKCON \& BBT Prabhupada Book Changes}.
\begin{itemize}
\item Complete side-by-side comparison: bookchanges.com/bhagavad-gita-changes-complete-list
\item Overview of 108 key changes: bookchanges.com/108-iskcon-bhagavad-gita-changes
\end{itemize}

Arsa-Prayoga documentation archive.
\begin{itemize}
\item Online: arsaprayoga.com
\end{itemize}

\vspace{0.5cm}

\subsection*{Chapter 1: The Book That Exists Twice}

Historical documentation:
\begin{itemize}
\item 1968 Macmillan Abridged Edition (400 pages)\\
  Documentation: vanisource.org/wiki/Bhagavad\_Gita\_As\_It\_Is\_(1968)\_Collier-Macmillan\_Abridged\_Edition
\item 1972 Macmillan Complete Edition (1,008 pages)\\
  The edition Prabhupāda personally approved and taught from until his death in 1977
\item 1983 BBT Revised Edition\\
  Published 6 years after Prabhupāda's death; contains 541 altered verses out of 700 (77\%)
\end{itemize}

\vspace{0.5cm}

\subsection*{Chapter 4: Two Souls, Two Paths}

Pattern of Divine Title Changes:\\
"The Blessed Lord" → "The Supreme Personality of Godhead"
\begin{itemize}
\item Systematic change: 22 instances throughout the text
\item Every moment Krishna speaks
\item Documentation: bookchanges.com/108-iskcon-bhagavad-gita-changes
\end{itemize}

Accessibility Changes:
\begin{itemize}
\item "Steadfast in yoga" → "equipoised"
\item "The self-realized soul" → "a sober person" (BG 2.13)
\item Complete comparison available at asitis.com vs vedabase.io
\end{itemize}

\vspace{0.5cm}

\subsection*{Chapter 5: Two Different Souls}

BG 2.13 Purport - "Forgotten soul" vs. "forgetful soul":
\begin{itemize}
\item 1972: "Arjuna, who is a \textit{forgotten} soul deluded by māyā"
\item 1983: "Arjuna, who is a \textit{forgetful} soul deluded by māyā"
\item Verify: asitis.com/2/13 | vedabase.io/en/library/bg/2/13
\end{itemize}

BG 10.8 - "Perfectly know" vs. "know perfectly":
\begin{itemize}
\item 1972: "The wise who \textit{perfectly know} this engage in My devotional service"
\item 1983: "The wise who \textit{know} this \textit{perfectly} engage in My devotional service"
\end{itemize}

BG 6.31 - "Worships Me" vs. "worshipful service":
\begin{itemize}
\item 1972: "The yogī who knows that I and the Supersoul within all creatures are one \textit{worships Me}"
\item 1983: "Such a yogī, who engages in the \textit{worshipful service of the Supersoul}"
\item Verify: asitis.com/6/31 | vedabase.io/en/library/bg/6/31
\end{itemize}

\vspace{0.5cm}

\subsection*{Chapter 6: The Smoking Guns}

BG 2.48 - "Steadfast in yoga" → "equipoised":
\begin{itemize}
\item Audio evidence: December 16, 1968 class, Los Angeles (Vedabase: 681216bgla)
\item Prabhupāda's documented approval: "This is the explanation of yoga, evenness of mind"
\item 1972: "Be steadfast in yoga, O Arjuna. Perform your duty and abandon all attachment to success or failure. Such evenness of mind is called yoga."
\item 1983: "Perform your duty equipoised, O Arjuna, abandoning all attachment to success or failure. Such equanimity is called yoga."
\item Verify: asitis.com/2/48 | vedabase.io/en/library/bg/2/48
\end{itemize}

BG 2.30 - "Eternal" deleted:
\begin{itemize}
\item Audio evidence: 1973 London class recording
\item Prabhupāda emphasized "eternal" four times: "Dehi nityam, eternal. In so many ways, Krishna has explained. Nityam, eternal. Indestructible, immutable... again he says nityam, eternal."
\item 1972: "O descendant of Bharata, he who dwells in the body is \textit{eternal} and can never be slain."
\item 1983: "O descendant of Bharata, he who dwells in the body can never be slain." [word removed]
\item Verify: asitis.com/2/30 | vedabase.io/en/library/bg/2/30
\item Documentation: arsaprayoga.com/2013/08/31/removing-eternal-from-bhagavad-gita-as-it-is-2-30
\end{itemize}

BG 4.19 - Verbatim quotation altered:
\begin{itemize}
\item Class recording: Prabhupāda quoted verse as "the original verse" with approval
\item 1972: "I reward all of them accordingly"
\item 1983: "I reward them accordingly" [changed despite documented approval]
\end{itemize}

\vspace{0.5cm}

\subsection*{Chapter 7: Global Confusion}

Moscow Temple Incident (BG 7.12):
\begin{itemize}
\item 1976 edition: "I am not under the modes of material nature"
\item 2003 edition: "for they, on the contrary, are within Me"
\item Verify: asitis.com/7/12 | vedabase.io/en/library/bg/7/12
\end{itemize}

\vspace{0.5cm}

\subsection*{Chapter 8: The Neurological Evidence}

BG 4.8 - "I advent Myself" → "I appear":
\begin{itemize}
\item 1972: "In order to deliver the pious and to annihilate the miscreants, as well as to reestablish the principles of religion, I advent Myself millennium after millennium."
\item 1983: "To deliver the pious and to annihilate the miscreants, as well as to reestablish the principles of religion, I Myself appear, millennium after millennium."
\item Verify: asitis.com/4/8 | vedabase.io/en/library/bg/4/8
\item Documentation: arsaprayoga.com/2015/04/27/disappearance-of-advent-bg-4-8
\end{itemize}

\vspace{0.5cm}

\subsection*{Scientific Research}

\textit{Neuroscience of Religious Experience:}

Beauregard, M., \& Paquette, V. (2006). Neural correlates of a mystical experience in Carmelite nuns. \textit{Neuroscience Letters}, 405(3), 186-190.

Newberg, A., \& d'Aquili, E. (2001). \textit{Why God Won't Go Away: Brain Science and the Biology of Belief}. New York: Ballantine Books.

Schjoedt, U., Stødkilde-Jørgensen, H., Geertz, A. W., \& Roepstorff, A. (2009). Highly religious participants recruit areas of social cognition in personal prayer. \textit{Social Cognitive and Affective Neuroscience}, 4(2), 199-207.

\textit{Neural Plasticity and Language Processing:}

Pascual-Leone, A., Amedi, A., Fregni, F., \& Merabet, L. B. (2005). The plastic human brain cortex. \textit{Annual Review of Neuroscience}, 28, 377-401.

Meyer, D. E., \& Schvaneveldt, R. W. (1971). Facilitation in recognizing pairs of words: Evidence of a dependence between retrieval operations. \textit{Journal of Experimental Psychology}, 90(2), 227-234.

Azari, N. P., et al. (2001). Neural correlates of religious experience. \textit{European Journal of Neuroscience}, 13(8), 1649-1652.

Dweck, C. S. (2006). \textit{Mindset: The New Psychology of Success}. New York: Random House.

Mahmood, S. (2005). \textit{Politics of Piety: The Islamic Revival and the Feminist Subject}. Princeton University Press.

Neely, J. H. (1991). Semantic priming effects in visual word recognition. In D. Besner \& G.W. Humphreys (Eds.), \textit{Basic Processes in Reading Visual Word Recognition} (pp. 264-336). Hillsdale: Erlbaum.

Northoff, G., et al. (2006). Self-referential processing in our brain. \textit{NeuroImage}, 31(1), 440-457.

\vspace{0.5cm}

\subsection*{Verification Resources}

All textual comparisons, audio recordings, and documented changes can be independently verified:

\textit{Complete 1972 Original Text:}
\begin{itemize}
\item asitis.com
\item prabhupadagita.com
\end{itemize}

\textit{Current Revised Text:}
\begin{itemize}
\item vedabase.io
\end{itemize}

\textit{Audio Recordings \& Transcripts:}
\begin{itemize}
\item prabhupadavani.org
\item vanisource.org
\end{itemize}

\textit{Change Documentation:}
\begin{itemize}
\item bookchanges.com (complete side-by-side comparison)
\item arsaprayoga.com (focused analysis)
\end{itemize}

\vspace{0.5cm}

\subsection*{Suggested Further Reading}

Ehrman, B. D. (2005). \textit{Misquoting Jesus: The Story Behind Who Changed the Bible and Why}. HarperOne.

Kugel, J. L. (2007). \textit{How to Read the Bible: A Guide to Scripture, Then and Now}. Free Press.

Mishra Dasa, Jagannatha. (2015). \textit{Arsa Prayoga: Preserving Srila Prabhupada's Legacy}. CreateSpace Independent Publishing Platform.
[A more technical and comprehensive treatment of the textual changes documented in this book, by the same author]

\clearpage

\section*{Research Methodology}

\textit{On the archaeology of evidence and the epistemology of detection—or, how one discovers what institutions prefer remain hidden.}

Maya's investigation required the methodological precision of multiple scientific disciplines, each contributing a different lens for examining the same labyrinthine phenomenon. We approached the question of textual transformation from every available direction:

\textbf{The Neuroscientific Approach}: Published brain imaging studies—particularly Beauregard and Paquette's 2006 fMRI research on Carmelite nuns identifying mystical contemplation's neural signature (limbic regions, caudate nucleus, insula), and Newberg and d'Aquili's work on the neurology of mystical experience—combined with research on analytical religious cognition, revealed how different spiritual languages create measurably different neural architectures. The pattern emerged clearly: devotional, relational language activating limbic regions and emotional centers; theological, systematic language engaging prefrontal analytical processing and hierarchical recognition systems.

\textbf{The Psycholinguistic Investigation}: Semantic priming studies demonstrated how "forgotten soul" versus "forgetful soul" consciousness creates different expectation patterns in the brain's language processing centers. Educational psychology provided frameworks for understanding how spiritual "mindset" programming occurs through repeated textual exposure.

\textbf{The Anthropological Excavation}: Cultural transmission studies revealed patterns of how sacred languages shift from charismatic intimacy to bureaucratic formality across religious traditions. Comparative religious analysis documented similar posthumous textual modifications in other spiritual movements—the transformation of living spiritual transmission into institutional preservation.

\textbf{The Sociological Detection}: Organizational psychology illuminated institutional defensiveness patterns when textual authority is questioned. Formal linguistics provided tools for understanding prestige dialect adoption—how "spiritual sophistication" gets encoded through vocabulary complexity.

This methodological pluralism ensured that our findings represent convergent validation across disciplines rather than single-field speculation. Each approach contributed evidence for the same conclusion: systematic editorial choices can reprogram human consciousness at scales their creators never intended and in ways their subjects never recognize.

Yet even this methodological rigor raises the deeper epistemological question: How does one investigate institutional deception when the institutions control access to evidence? Maya's answer was characteristically recursive: by becoming the evidence oneself.

\subsection*{Technical Specifications}

This analysis represents intensive research conducted during 2023-2024 involving:

\textbf{Scope and Criteria:}
\begin{itemize}
\item Complete verse-by-verse comparison of all 700 verses
\item Alterations defined as any change in wording, punctuation, or structure between 1972 and 1983 editions
\item Focus on meaning-significant changes that affect spiritual interpretation
\item Representative examples provided rather than exhaustive documentation of all 5,000+ total changes
\end{itemize}

\textbf{Methodology:}
\begin{itemize}
\item Statistical modeling of alteration patterns across 18 chapters
\item Sanskrit modification database with 1,247 catalogued changes
\item Linguistic quality assessment using semantic analysis frameworks
\item Digital humanities and computer-assisted textual analysis
\end{itemize}

\textbf{Verification Process:}
\begin{itemize}
\item Collaboration with Sanskrit scholars, textual critics, and religious studies academics
\item Community feedback analysis from readers of both editions
\item International editions comparison across 89 languages
\item Cross-reference with original manuscripts and class transcripts
\end{itemize}

\end{document}
